% registry: KHALIJ
\methodhead{KHALIJ}{Khalij, Pagnacco and Troian (2017): improved Gough / Nishihara--Kawamoto criterion}\label{sheet:KHALIJ}

\block{Source} L.~Khalij, E.~Pagnacco, R.~Troian, \emph{Int.\ J.\ Fatigue} 99 (2017) 137--150
\cite{khalij2017}, Eqs.~(12), (22), (23), (26); Tables~1 and~6.

\block{Idea and scope} A fatigue-limit criterion for zero-mean sinusoidal plane stress that
weights an equivalent shear stress amplitude against the amplitude of the first invariant. The
weight depends on the limit ratio $r=\tm/\fm$, so one expression moves from a von Mises ellipse
(mild steels, $r\approx 1/\sqrt3$) to a Rankine-like envelope (cast irons, $r\to1$). Proposed for
metals with $0.5\lesssim r\le 1$, in-phase and out-of-phase loading.

\block{Inputs}
\begin{center}\small
\begin{tabular}{@{}llp{100mm}@{}}
\toprule
Symbol & Unit & Meaning\\\midrule
$s_{xx,a}$, $s_{yy,a}$, $s_{xy,a}$ & MPa & amplitudes of the plane-stress components\\
$\phi_{yy}$ & rad & phase shift between the two normal stresses $s_{xx}$ and $s_{yy}$\\
$s_{-1}=\fm$, $\tm$ & MPa & fully reversed axial (or bending) and torsion limits\\
$r=\tm/\fm$ & -- & limit ratio\\
$\alpha(r)$ & -- & weight of the shear term, Eq.~\eqref{eq:khalij-alpha}\\
$J_{2,a}$, $I_{1,a}$ & MPa$^2$, MPa & PHPS shear amplitude squared and first-invariant amplitude\\
$\mathfrak D$ & -- & damage; $\mathfrak D\le1$ means no failure\\
\bottomrule
\end{tabular}
\end{center}

\block{Formulation} With the PHPS amplitude (Eq.~12), which depends on the phase of the normal
stresses only, and the amplitude of the first invariant (Eq.~22),
\begin{gather}
J_{2,a} = \tfrac13\left(s_{xx,a}^2 + s_{yy,a}^2 - s_{xx,a}s_{yy,a}\cos\phi_{yy}\right) + s_{xy,a}^2 ,\qquad
I_{1,a} = \sqrt{s_{xx,a}^2 + s_{yy,a}^2 + 2\,s_{xx,a}s_{yy,a}\cos\phi_{yy}} ,\\
\mathfrak D = \alpha(r)\left[\Bigl(\frac{s_{xx,a}}{\fm}\Bigr)^2 + \Bigl(\frac{s_{yy,a}}{\fm}\Bigr)^2
 - \frac{s_{xx,a}}{\fm}\frac{s_{yy,a}}{\fm}\cos\phi_{yy}\right]
 + \Bigl(\frac{s_{xy,a}}{\tm}\Bigr)^2 + \bigl(1-\alpha(r)\bigr)\frac{I_{1,a}}{\fm} \le 1 ,
 \label{eq:khalij}\\
\alpha(r) = \frac{1}{3r^2} . \label{eq:khalij-alpha}
\end{gather}
Eq.~\eqref{eq:khalij-alpha} is the authors' setting ``OP'': it makes the criterion independent of
the frame for pure shear (Eqs.~24--26). With it the first two groups of Eq.~\eqref{eq:khalij} are
$J_{2,a}/\tm^2$, so
\begin{equation}
\mathfrak D = \frac{J_{2,a}}{\tm^2} + \Bigl(1 - \frac{1}{3r^2}\Bigr)\frac{I_{1,a}}{\fm} .
\label{eq:khalij-invariant}
\end{equation}
Both limits give $\mathfrak D=1$ exactly. For $r<1/\sqrt3$ the weight of $I_{1,a}$ is negative.

\block{Constants and calibration} None fitted: $\alpha$ follows from $\fm$ and $\tm$.

\block{Life} The paper gives none; finite lives use the master curve on $\mathfrak D$
(Section~\ref{sec:mastercurve}), a repository extension. $\mathfrak D$ is not homogeneous in the
load (linear $I_{1,a}$ term).

\block{Repository choices}
\begin{itemize}
\item Eq.~\eqref{eq:khalij-invariant} is evaluated with $J_{2,a}$ and $I_{1,a}$ of
  Eq.~\eqref{eq:invariants}; in plane stress they equal the paper's expressions.
\item Zero mean stress only. The paper's mean-stress extension (Eqs.~29--31) is not implemented:
  its term $\gamma_2(r)\,s_{ms}/s_{-1}^2$ is not dimensionless ($s_{ms}$ is linear in the mean
  stresses), Eq.~(30) for $\gamma_2$ follows from a purely quadratic base criterion, and no reading
  reproduces Table~7. A case with a mean stress, or with an out-of-plane component, is out of scope.
\item $r$ is computed from the material limits; Table~6 uses this unrounded ratio, not the
  two-digit $r$ printed beside it.
\end{itemize}

\block{Worked example} Table~6, CrV-V ($\fm=530$, $\tm=309\MPa$, $r=0.583$, $\alpha=0.981$), biaxial
limit $s_{xx,a}=0.24\fm=127.2$, $s_{yy,a}=-0.84\fm=-445.2\MPa$ (in phase, opposite signs):
$J_{2,a}=\tfrac13(127.2^2+445.2^2+127.2\cdot445.2)=90\,337\MPa^2$, $I_{1,a}=318.0\MPa$,
$\mathfrak D = 90\,337/309^2 + (1-0.981)\,318.0/530 = 0.946+0.012 = 0.958$ (printed 0.95).

\block{Verification} Table~6, OP column: 51 of 53 limits to one unit of the last printed digit;
two 3\%C cast-iron rows (1.035 and 1.063 against 1.06 and 1.01) are unexplained.
